\honest{Local Validity as a Key to Sanity and Civilization}{Eliezer Yudkowsky}{2018}

I claim that three things are alike. A valid step in a proof never turns a truth into a falsehood. A good reasoner judges an argument apart from whether they like its conclusion. And civilization needs coordinated steps away from bad equilibria, which need general rules that may punish people we like and reward people we dislike. When people stop believing the rules are applied fairly, they go back to the bad equilibrium and civilization falls. \nb{The title calls this a key to sanity and civilization; the post's own summary calls it a similarity.}

If $x = y$, then $2x = 2y$; the step is valid even in a model where $x \neq y$, since it adds no new problem. But $xy = xz$ does not give $y = z$: with $x = 0$, $y = 4$, $z = 17$, a true ``0 = 0'' yields a false ``4 = 17.'' That is the last proof in the essay. You cannot have the idea of a proof without approving of single steps apart from liking the conclusion, and once your grasp of valid steps is stronger than your faith in your intuitive verdicts, you can be convinced of counterintuitive truths.

Then I extend this to all arguments. Bad arguments for true conclusions are easy to find, ``because there are tiny elves that whisper them to people.'' An ``intellectually strong mind'' judges an argument apart from its conclusion. Someone who doubts the intelligence explosion can still frown at the argument that it is impossible because hypercomputation is impossible; someone who accepts global warming still winces at a scorching day offered as evidence. One of my examples of a terrible argument is that there is ``no such thing as intelligence'' because of the no-free-lunch theorem, and it links to my reply to François Chollet. \nb{Chollet wrote that there is no such thing as ``general'' intelligence. The post drops the word ``general.''}

Some people reason that the hot day was part of a pattern of record highs, so it counts; others reason that if they would roll their eyes at a cold day offered against warming, a hot day cannot count for it. I would pay a 5\% premium to buy a used car from the second. The first will court-martial an allied argument if they must but favor allies when they can. The second, like a mathematician, praises valid steps on either side without feeling disloyal.

Good arguments for false conclusions are much rarer, since strong evidence is by definition rarely found for a false conclusion. My example is Lord Kelvin, whose arguments that the Earth was young ``all failed simultaneously'' because his era did not know about nuclear physics. \nb{For his cooling-Earth calculation, that is wrong. The geophysicists England, Molnar and Richter have shown that adding radioactivity to it does not change his answer. His error was to assume that heat leaves the Earth's interior only by conduction, and John Perry pointed it out in 1895, before radioactivity was known. His argument from the age of the Sun did fail for want of nuclear physics.}

The same holds in law. A potential juror in the trial of Martin Shkreli said, ``I can be fair to one side but not the other,'' and was rejected during jury selection. I call this ``possibly a harbinger of the collapse of civilization,'' and note that much of the code of fairness is not written down anywhere. \nb{The evidence for this harbinger is one prospective juror.}

Next, Democrats debating whether to push Al Franken out of the Senate, worried that the Republicans would not police their own side. I set aside whether Franken deserved it. Some replied that this was horrifyingly cynical, as if misconduct were punishable only in Republicans; others that the alternative was handing the Senate to the Republicans. I think the knot comes from not separating game theory from goodness. I have heard less of that worry since ``some upstanding Republican voters in Alabama stayed home'' and elected Doug Jones. \nb{About a week before that vote, the Republican president had endorsed the accused Republican candidate, and the party had rejoined his campaign. So the policing I credit came from voters only, not from the party.}

Ideally law is the part of morality wise to enforce with guns, like ``Don't kill people''; laws against marijuana show it falling short. From that view, sparing one's own senators looks like giving up. But law, like money with its two functions, is also game theory among people who may not share your morals: it lets them move from mutual defection at (2, 2) to enforced cooperation at (2.9, 2.9), each paying 0.1 for enforcement.

On that view everything rests on impartiality, on not caring whose ox is gored; if the law systematically punishes your defection and not the other side's, it pays you to blow up the equilibrium. So it is coherent to call the conduct bad whoever does it and still refuse to punish only one's own side, though Alabama shows the danger of assuming the other side will not cooperate. Some people say they would gladly take my stuff if the law did not stop them; I am not sure how seriously to take that. Enforcement is ``simply not effective enough to count for the vast majority of human cooperation,'' and civilization ``is free-riding a whole lot on innate altruism.'' \nb{If so, it is unclear how unfair enforcement alone brings civilization down, and the essay does not say.}

Human law depends on rules that feel simple and general, because of our cognitive limits; two superintelligences could agree on a compromise with complicated boundaries right up to the Pareto frontier. The laws I mean are not the modern morass of regulations but those a small town enforced in 1820, or that cops enforce against other cops. Hunter-gatherers lack 100,000 pages of law not from wisdom about loopholes and unintended consequences, but for lack of time and paper. Spoken law comes out as short universal sentences, like ``If you kill somebody who wasn't attacking you first, we'll exile you.''

Writing the law down let everyone know what it was and be safe if they kept it. ``Robert Heinlein called it the most important moment in political history.'' I give no source for this. I doubt the Code of Hammurabi was universally enforced, and suspect that most real law was never written and much that was written was never real law. Then people wrote down too much, and today almost nothing that serves the true function of law can be written down. The criminal system is too slow and unreliable for any sane victim of sexual assault to prefer it to the media or the whispernet; civil law is a bludgeoning contest between those who can afford lawyers; and in poor American neighborhoods the police will not let residents form their own order.

What remains of law must be known without being written and applied by the community without professional judges; at least the elders must seem to appeal to the law and not to self-interest, and it breaks down if people believe deciders look at whose ox was gored. It feels simple the way walking feels easy, though it is not computationally simple. So human law is what people at least ``believe to be a set of simple rules that can be locally checked to test okay behavior.'' \nb{The one case in the essay of police enforcing the law unequally, seizing the cars of black people but not white people who use marijuana, appears in a clause beginning ``regardless of whether.''}

Following such rules can feel like losing: you give up the (5, 0) payoff for (3, 3), and the law may punish an ally. You can coherently think one enforcement makes the world worse while valuing the law, if you see its game-theoretic function. So long as the rules move everyone to a better equilibrium and fall on everyone equally, people can step back and weigh the law's interest against their own. That is called justice, fairness, impartiality.

My model of impartial law is ``a story I couldn't seem to Google'': a Chinese general whose son took a straw hat from a peasant, so the general sentenced his son to death, in tears. \nb{The story that matches is Lü Meng's, in the \textsc{Records of the Three Kingdoms}, and the man executed was a soldier from Lü Meng's home region.} My other model is a judge in H. Beam Piper's 1962 science fiction novel \textsc{Little Fuzzy}, who is ``just pro-law'': if the planet has native sapient inhabitants, the company's charter is illegal, and ``Frederic Pendarvis' religion is the law, and he is its priest.'' In 1962 no one in the story thinks him naive or strange. Without people like Pendarvis, and without appreciating them even when they rule against you, your tribe ends sooner or later. I doubt the United States will fall into anarchy this way, but the point applies down to bargains between two or three people.

Saying you can be fair to one side but not the other treats fairness as a favor to friends, and shows that even the instinctive sense of law as game theory is being lost. People are losing the ability to appreciate such stories because social media shows them only the other side defecting and only their own side's views, ``Or maybe'' because of chemicals leached from plastic food containers.

People appreciate impartial rules more natively than valid proof steps, so many start epistemology by seeing the rules of argument as law and fallacies as crimes; the unusually healthy reject bad allied arguments with the feeling of an impartial judge. Then I concede, in two sentences, that ``there is no game theory and no morality to the true way of the map that reflects the territory.'' \nb{So valid proof and fair law work in different ways. What they share is a feeling of following rules.}

I end with a guess: you are better off buying a used car from a random mathematician than from a random non-mathematician, because mathematicians' ``sense of Law was strong enough to be appropriated for proofs.'' I call this ``a guess and hypothesis'' and say I would like to see a study, though not one using self-reports of crime, which may measure honesty times criminality. I have no grand agenda; I only wish more of the most basic rules of thinking were written down.
